% P11 paper module: R18 off-diagonal near-null block and finite source-block firewall.

\subsection{Off-diagonal near-null suppression and the finite source-block firewall}
\label{subsec:o3q-offdiagonal-nearnull}

Keep the fixed smooth odd pair $f_0,f_m\in C_c^\infty((-R,R))$ from
Proposition~\ref{prop:o3n-nearnull}, with first integral-jet orders $0$ and $m>0$, and
\[
 z_U=f_m-\frac{\ell_U(f_m)}{\ell_U(f_0)}f_0,
 \qquad \ell_U(z_U)=0.
\]
Recall from Theorem~\ref{thm:o3p-nearnull-vanish} that
\[
 D_U(z_U,z_U)=\sigma_U(J_{R,U}z_U,J_{R,U}z_U)\to0.
\]
The mixed entry on the pair $(f_0,z_U)$ is then forced below the natural square-root
boundary scale.

\begin{theorem}[Off-diagonal TC1 suppression]\label{thm:o3q-mixed-suppression}
One has
\[
 \boxed{
 \rho_U(f_0,z_U)=0,
 \qquad
 \sigma_U(J_{R,U}f_0,J_{R,U}z_U)
 =o\!\left(\frac{e^{U/2}}{U}\right).
 }
\tag{O3Q.1}\label{eq:o3q-mixed-suppression}
\]
\end{theorem}

\begin{proof}
Since $\ell_U(z_U)=0$,
\[
 \rho_U(f_0,z_U)
 =\frac{\ell_U(f_0)\overline{\ell_U(z_U)}}{d_U}=0,
\]
so
\[
 \sigma_U(Jf_0,Jz_U)=D_U(f_0,z_U).
\]
The positive Gram representation~\eqref{eq:mj-remainder-gram} gives
\[
 |D_U(f_0,z_U)|^2
 \le D_U(f_0,f_0)D_U(z_U,z_U).
\]
For the fixed first-jet-zero vector $f_0$, equation~\eqref{eq:mj-Ddiag} gives
\[
 D_U(f_0,f_0)=o(e^U/U^2),
\]
while Theorem~\ref{thm:o3p-nearnull-vanish} gives
$D_U(z_U,z_U)\to0$.  Hence
\[
 |D_U(f_0,z_U)|^2=o(e^U/U^2),
\]
which proves~\eqref{eq:o3q-mixed-suppression}.
\end{proof}

The full graph form has a stronger geometric consequence because R17 leaves a fixed
positive Gamma floor on $z_U$.

\begin{corollary}[Asymptotic full-Gram orthogonality]\label{cor:o3q-full-angle}
Put
\[
 L_U:=\frac{e^U}{U^2}.
\]
There is $a_0>0$ such that
\[
 q_U^X(J_{R,U}f_0)=a_0L_U(1+o(1)),
\tag{O3Q.2}\label{eq:o3q-f0-diag}
\]
while
\[
 q_U^X(J_{R,U}z_U)
 \longrightarrow
 \gamma_m:=\mathfrak c_{\Gamma,R}[f_m]>0.
\tag{O3Q.3}\label{eq:o3q-z-diag}
\]
Moreover
\[
 q_U^X(J_{R,U}f_0,J_{R,U}z_U)=o(L_U^{1/2}),
\tag{O3Q.4}\label{eq:o3q-full-mixed}
\]
and therefore
\[
 \boxed{
 \frac{q_U^X(J_{R,U}f_0,J_{R,U}z_U)}
 {q_U^X(J_{R,U}f_0)^{1/2}
  q_U^X(J_{R,U}z_U)^{1/2}}
 \longrightarrow0.
 }
\tag{O3Q.5}\label{eq:o3q-full-angle}
\]
Equivalently, for the coordinate Gram matrix
\[
 M_U:=
 \begin{pmatrix}
 q_U^X(Jf_0) & q_U^X(Jf_0,Jz_U)\\
 q_U^X(Jz_U,Jf_0) & q_U^X(Jz_U)
 \end{pmatrix}
\]
and
\[
 S_U:=\operatorname{diag}
 \bigl(q_U^X(Jf_0)^{-1/2},q_U^X(Jz_U)^{-1/2}\bigr),
\]
one has
\[
 \boxed{S_UM_US_U\to I_2,}
\tag{O3Q.6}\label{eq:o3q-correlation-matrix}
\]
and
\[
 \boxed{\det M_U\sim a_0\gamma_mL_U.}
\tag{O3Q.7}\label{eq:o3q-det}
\]
\end{corollary}

\begin{proof}
Theorem~\ref{thm:odd} with first jet $0$ gives the Schur asymptotic in
\eqref{eq:o3q-f0-diag}; the fixed Gamma term is $O(1)$ and hence negligible on the
scale $L_U$.  Equation~\eqref{eq:o3q-z-diag} is
Corollary~\ref{cor:o3p-gamma-limit}.

For the mixed form, terminal Gamma compatibility gives
\[
 q_U^X(Jf_0,Jz_U)
 =\mathfrak c_{\Gamma,R}[f_0,z_U]
  +\sigma_U(Jf_0,Jz_U).
\]
Since $z_U\to f_m$ in the fixed smooth source space,
$\mathfrak c_{\Gamma,R}[f_0,z_U]=O(1)$.  Theorem~\ref{thm:o3q-mixed-suppression}
therefore proves~\eqref{eq:o3q-full-mixed}.  Dividing by the two diagonal asymptotics
gives~\eqref{eq:o3q-full-angle}, and the matrix statements follow immediately.
\end{proof}

The preceding two-vector statement does not create a source-level gauge diagnostic.
In fact the cocycle trivializes every finite block made only from canonically nested
source vectors.

\begin{proposition}[Complete finite source-block cocycle invariance]
\label{prop:o3q-finite-block-invariance}
Let $0<R<S<T$ and choose any finite family
$e_1,\dots,e_n\in\KXR{R}$.  Put
\[
 e_i^S:=J_{R,S}e_i.
\]
Then for every terminal $T>S$ and every $i,j$,
\[
 \boxed{
 \langle G_{S,T}e_i^S,e_j^S\rangle_{X,S}
 =
 \langle G_{R,T}e_i,e_j\rangle_{X,R}.
 }
\tag{O3Q.8}\label{eq:o3q-block-invariance}
\]
Consequently, for fixed $T_0$ and future terminal $U$, the complete finite Gram pair
\[
 \left(
 [\langle G_{R,T_0}e_i,e_j\rangle],
 [\langle G_{R,U}e_i,e_j\rangle]
 \right)
\]
is identical, under the canonical coordinate identification, to the corresponding
pair after transport to level $S$.

The same assertion holds pointwise for $U$-dependent finite families, provided their
$S$-level representatives are the canonical images of the $R$-level vectors.
\end{proposition}

\begin{proof}
By the graph cocycle,
\[
 J_{S,T}e_i^S
 =J_{S,T}J_{R,S}e_i
 =J_{R,T}e_i.
\]
Since $G_{X,T}=J_{X,T}^*J_{X,T}$ with the graph-Hilbert adjoint,
\[
 \langle G_{S,T}e_i^S,e_j^S\rangle_{X,S}
 =\langle J_{S,T}e_i^S,J_{S,T}e_j^S\rangle_{X,T}
 =\langle J_{R,T}e_i,J_{R,T}e_j\rangle_{X,T},
\]
which is~\eqref{eq:o3q-block-invariance}.  Applying it at $T_0$ and at $U$ proves the
Gram-pair statement.  The proof is pointwise in the chosen vectors and therefore also
applies to a $U$-dependent family.
\end{proof}

\begin{corollary}[The $(f_0,z_U)$ mixed scalar is not an $R/S$ gauge diagnostic]
\label{cor:o3q-mixed-no-gauge}
Let $0<R<S$ and regard the pair $(f_0,z_U)$ at level $S$ by canonical zero extension.
Then its complete terminal $(T_0,U)$ Gram pair, including every off-diagonal entry, is
identical at the two source levels.  In particular no difference of
\[
 q_U^X(Jf_0,Jz_U),
 \qquad
 \sigma_U(Jf_0,Jz_U),
\]
or of any finite generalized Gram quantity built solely from the canonically nested
pair can measure the $R$--$S$ polar-gauge defect.
\end{corollary}

\begin{proof}
Apply Proposition~\ref{prop:o3q-finite-block-invariance} to the two-vector family.  The
Schur statement follows as well because the Gamma form is terminal-compatible, so
subtracting the identical Gamma matrix from the identical full terminal Gram matrix
leaves the identical Schur matrix.
\end{proof}

\begin{remark}[R18 firewall: internal blocks versus full square roots]
\label{rem:o3q-firewall}
Corollary~\ref{cor:o3q-full-angle} resolves the normalized orientation of the
\emph{internal two-vector future Gram block}: after exact boundary-mode cancellation the
large $f_0$ direction and the Gamma-supported $z_U$ direction are asymptotically
orthogonal.  Proposition~\ref{prop:o3q-finite-block-invariance} shows that this complete
finite form pair is already source-compatible under $R\to S$.

This does not determine the full polar gauges.  Compression and positive functional
calculus do not commute in general:
\[
 P_EA^{1/2}P_E\ne(P_EAP_E)^{1/2}.
\]
The canonical-inclusion model of Subsection~\ref{subsec:o3m-polar-gauge} already shows
that perfect internal modulus data need not determine the true future transport.
Therefore the remaining gauge information must involve the coupling of the nested
source image with its complement before the full square root/polar factor is taken.

A genuine next target is consequently an off-block square-root quantity in the
$T_0$-orthogonal splitting
\[
 \Ran W\oplus(\Ran W)^\perp,
\]
for example the relation between the already nontrivial block
\[
 (I-WW^*)A_SW
\]
and its square-root analogue
\[
 (I-WW^*)A_S^{1/2}W,
\]
or directly the polar leakage $(I-WW^*)U_SW$.
No conclusion about $\Gamma_U\to I$, $K_{R,S}^{T_0,U}\to I$, strong terminal
transport, Object X, Seal, or RH is drawn here.
\end{remark}
